\documentclass[10pt]{article}
\usepackage[a4paper,margin=1.8cm]{geometry}
\usepackage{booktabs}
\usepackage{array}
\usepackage{longtable}
\usepackage{caption}
\usepackage{float}
\setlength{\parskip}{3pt}
\setlength{\tabcolsep}{3pt}
\renewcommand{\arraystretch}{1.08}
\title{Supplementary Evidence for\\Explainable Early Warning for Next-Week Abnormal Reported 311 Demand}
\author{}
\date{}
\begin{document}
\maketitle
\noindent This supplement is generated from the repository's current stored artifacts. The final comparison uses the T2 event: a shifted eight-week abnormality and next-week count at least three. Unless stated otherwise, the final protocol trains through 2023, selects calibration and thresholds on 2024, and evaluates unchanged on held-out 2025 rows.

\section*{S1. Historical-volume performance}
Train-derived cut points of the shifted eight-week historical mean define the volume groups; no test label is used to form the groups. Repeated zero-valued cut points collapse the nominal deciles to nine observable bins. All rows below sum to the full held-out 2025 T2 population.
\begin{table}[H]
\centering\scriptsize
\begin{tabular}{@{}lrrrrrrrr@{}}
\toprule
Volume group & Rows & Prevalence & PR-AUC & Precision & Recall & F1 & P@5\% & Alert rate \\
\midrule
D1 & 22,986 & 1.0\% & 0.085 & 0.142 & 0.070 & 0.094 & 0.071 & 0.5\% \\\\
D2 & 11,784 & 11.1\% & 0.249 & 0.227 & 0.538 & 0.319 & 0.341 & 26.2\% \\\\
D3 & 10,915 & 14.1\% & 0.295 & 0.238 & 0.639 & 0.346 & 0.410 & 37.9\% \\\\
D4 & 11,920 & 14.8\% & 0.325 & 0.274 & 0.584 & 0.373 & 0.455 & 31.5\% \\\\
D5 & 12,453 & 14.4\% & 0.311 & 0.266 & 0.568 & 0.363 & 0.409 & 30.7\% \\\\
D6 & 12,491 & 12.6\% & 0.277 & 0.247 & 0.556 & 0.342 & 0.368 & 28.3\% \\\\
D7 & 11,129 & 12.7\% & 0.310 & 0.264 & 0.547 & 0.356 & 0.427 & 26.3\% \\\\
D8 & 11,432 & 13.3\% & 0.339 & 0.277 & 0.618 & 0.382 & 0.458 & 29.6\% \\\\
D9 & 17,506 & 13.9\% & 0.399 & 0.304 & 0.599 & 0.403 & 0.553 & 27.5\% \\\\
\bottomrule
\end{tabular}

\caption{Final Platt-calibrated LightGBM performance by train-derived historical-volume group on held-out 2025 T2 rows.}\label{tab:supp-volume}
\end{table}
\clearpage

\section*{S2. Complaint-category mapping and residual category}
\begin{table}[H]
\centering\small
\begin{tabular}{@{}lrrr@{}}
\toprule
Analysis category & Types & Requests & Share \\
\midrule
housing & 57 & 8,061,962 & 26.06\% \\\\
noise & 9 & 6,713,922 & 21.70\% \\\\
parking/traffic & 39 & 6,521,679 & 21.08\% \\\\
sanitation & 36 & 3,090,901 & 9.99\% \\\\
infrastructure & 28 & 2,673,125 & 8.64\% \\\\
public/safety & 38 & 1,597,023 & 5.16\% \\\\
water/sewer & 10 & 1,212,939 & 3.92\% \\\\
environment & 17 & 626,930 & 2.03\% \\\\
other & 86 & 441,648 & 1.43\% \\\\
\bottomrule
\end{tabular}

\caption{Deterministic mapping of original 311 complaint types to the nine analysis categories.}\label{tab:supp-mapping}
\end{table}
\begin{table}[H]
\centering\small
\begin{tabular}{@{}lrr@{}}
\toprule
Complaint type & Requests & Share of other \\
\midrule
Encampment & 191,408 & 43.34\% \\\\
Drug Activity & 93,873 & 21.26\% \\\\
Animal in a Park & 29,914 & 6.77\% \\\\
Vending & 27,575 & 6.24\% \\\\
Electronics Waste & 26,487 & 6.00\% \\\\
Miscellaneous Categories & 13,752 & 3.11\% \\\\
Unleashed Dog & 11,718 & 2.65\% \\\\
DOF Property - Payment Issue & 11,414 & 2.58\% \\\\
DRIE & 4,023 & 0.91\% \\\\
Storm & 3,003 & 0.68\% \\\\
Posting Advertisement & 2,745 & 0.62\% \\\\
DPR Internal & 2,520 & 0.57\% \\\\
\bottomrule
\end{tabular}

\caption{Largest original complaint types in the residual \texttt{other} category.}\label{tab:supp-other}
\end{table}
\clearpage

\section*{S3. Same-target baseline comparison}
Every row uses the same held-out 2025 T2 population. Formula-threshold decisions apply the T2 threshold to predicted counts; the LightGBM threshold is selected on validation 2024.
\begin{table}[H]
\centering\scriptsize
\begin{tabular}{@{}llrrrrrr@{}}
\toprule
Method & Decision & Count MAE & PR-AUC & P@5\% & Precision & Recall & F1 \\
\midrule
Poisson GLM, no NTA FE & Formula threshold & 9.069 & 0.141 & 0.146 & 0.376 & 0.038 & 0.068 \\\\
Poisson GLM + NTA FE & Formula threshold & 9.070 & 0.141 & 0.146 & 0.375 & 0.038 & 0.068 \\\\
HGB Poisson count & Formula threshold & 8.021 & 0.153 & 0.176 & 0.513 & 0.103 & 0.171 \\\\
Hurdle HGB count & Formula threshold & 7.938 & 0.152 & 0.172 & 0.497 & 0.111 & 0.181 \\\\
No-shortcut LightGBM classifier & Validation threshold & -- & 0.317 & 0.418 & 0.263 & 0.574 & 0.361 \\\\
\bottomrule
\end{tabular}

\caption{Full baseline metrics on the shared T2 task.}\label{tab:supp-baselines}
\end{table}
\clearpage

\section*{S4. Bootstrap uncertainty and seed stability}
Baseline intervals use 250 NTA-category cluster resamples; the main LightGBM confidence intervals reported in the manuscript use 1,000. The reported estimate is the held-out 2025 value; this table does not turn distinct decision rules into a single superiority claim.
\begin{table}[H]
\centering\scriptsize
\begin{tabular}{@{}llrr@{}}
\toprule
Model & Metric & Estimate & Cluster-bootstrap 95\% CI \\
\midrule
No-shortcut LightGBM & PR-AUC & 0.317 & [0.308, 0.326] \\\\
No-shortcut LightGBM & P@5\% & 0.418 & [0.403, 0.432] \\\\
No-shortcut LightGBM & F1 & 0.361 & [0.355, 0.367] \\\\
Poisson GLM, no NTA FE & PR-AUC & 0.141 & [0.138, 0.144] \\\\
Poisson GLM, no NTA FE & P@5\% & 0.146 & [0.136, 0.155] \\\\
Poisson GLM, no NTA FE & F1 & 0.068 & [0.062, 0.076] \\\\
Poisson GLM + NTA FE & PR-AUC & 0.141 & [0.138, 0.144] \\\\
Poisson GLM + NTA FE & P@5\% & 0.146 & [0.135, 0.157] \\\\
Poisson GLM + NTA FE & F1 & 0.068 & [0.061, 0.075] \\\\
HGB Poisson count & PR-AUC & 0.153 & [0.149, 0.157] \\\\
HGB Poisson count & P@5\% & 0.176 & [0.165, 0.184] \\\\
HGB Poisson count & F1 & 0.171 & [0.162, 0.180] \\\\
Hurdle HGB count & PR-AUC & 0.152 & [0.148, 0.156] \\\\
Hurdle HGB count & P@5\% & 0.172 & [0.164, 0.183] \\\\
Hurdle HGB count & F1 & 0.181 & [0.171, 0.189] \\\\
\bottomrule
\end{tabular}

\caption{Cluster-bootstrap uncertainty for final-comparison rows.}\label{tab:supp-bootstrap}
\end{table}
\begin{table}[H]
\centering\small
\begin{tabular}{@{}lrrrr@{}}
\toprule
Configuration & Seeds & PR-AUC & P@5\% & F1 \\
\midrule
03 history calendar & 5 & 0.288 $\pm$ 0.003 & 0.385 $\pm$ 0.004 & 0.346 $\pm$ 0.003 \\\\
04 history calendar weather & 5 & 0.298 $\pm$ 0.003 & 0.390 $\pm$ 0.004 & 0.353 $\pm$ 0.002 \\\\
07 final no shortcut borough & 5 & 0.319 $\pm$ 0.002 & 0.419 $\pm$ 0.003 & 0.364 $\pm$ 0.002 \\\\
08 final no shortcut nta fe & 5 & 0.319 $\pm$ 0.001 & 0.421 $\pm$ 0.001 & 0.364 $\pm$ 0.001 \\\\
\bottomrule
\end{tabular}

\caption{Five-seed held-out stability for key ablation configurations. Values are mean $\pm$ standard deviation.}\label{tab:supp-seeds}
\end{table}
\clearpage

\section*{S5. Chronological, calibration, and target-sensitivity evidence}
\begin{table}[H]
\centering\small
\begin{tabular}{@{}lrrrrr@{}}
\toprule
Test year & Rows & Prevalence & PR-AUC & P@5\% & F1 \\
\midrule
2021 & 122,616 & 11.4\% & 0.312 & 0.427 & 0.375 \\\\
2022 & 122,616 & 10.9\% & 0.281 & 0.373 & 0.349 \\\\
2023 & 122,616 & 10.9\% & 0.281 & 0.381 & 0.340 \\\\
2024 & 124,974 & 11.0\% & 0.308 & 0.408 & 0.366 \\\\
2025 & 122,616 & 11.1\% & 0.317 & 0.418 & 0.361 \\\\
\bottomrule
\end{tabular}

\caption{Expanding-window rolling-origin test metrics for T2.}\label{tab:supp-rolling}
\end{table}
\begin{table}[H]
\centering\small
\begin{tabular}{@{}lrrrr@{}}
\toprule
Method & Brier & ECE & Log loss & PR-AUC \\
\midrule
Uncalibrated & 0.168 & 0.249 & 0.493 & 0.317 \\\\
Platt & 0.087 & 0.006 & 0.293 & 0.317 \\\\
Isotonic & 0.087 & 0.004 & 0.292 & 0.307 \\\\
\bottomrule
\end{tabular}

\caption{Final-style 2025 calibration diagnostics for T2.}\label{tab:supp-calibration}
\end{table}
\noindent The following is a labelled 2024--2025 target-definition diagnostic, not the same-task baseline comparison in Section S3.
\begin{table}[H]
\centering\small
\begin{tabular}{@{}lrrrrr@{}}
\toprule
Target & Rows & Prevalence & PR-AUC & P@5\% & F1 \\
\midrule
T0 & 247,590 & 12.8\% & 0.310 & 0.424 & 0.360 \\\\
T1 & 247,590 & 11.7\% & 0.313 & 0.418 & 0.364 \\\\
T2 & 247,590 & 11.0\% & 0.316 & 0.415 & 0.366 \\\\
T3 & 189,423 & 13.7\% & 0.324 & 0.441 & 0.371 \\\\
\bottomrule
\end{tabular}

\caption{Target-sensitivity diagnostic with prevalence and ranking metrics.}\label{tab:supp-targets}
\end{table}

\section*{S6. Count-model limitation}
Poisson GLM fits converged within the configured 500 iterations. A Negative Binomial GLM is not included because \texttt{statsmodels} was unavailable in the recorded environment and scikit-learn does not provide a Negative Binomial GLM. This is a documented tooling limitation, not an omitted result.
\end{document}
